% ------------------------------- Chapter 1 -------------------------------
\chapter{Introduction}

A turbomachine produces or absorbs energy through a change of enthalpy in a
continuously flowing stream, transferring work through a rotating shaft. In a
gas turbine the turbine extracts work from the hot combustion products by
expansion: the flow accelerates at the expense of pressure, and a large part of
the extracted power is used to drive the compressor through the shaft, the
remainder being the useful (net) output.

Designing a turbine stage couples thermodynamics, kinematics (velocity
triangles), annulus geometry and mechanical constraints~\cite{dixon}. The
procedure followed in this work is the classical three-step chain of preliminary
design:

\begin{itemize}
    \item \textbf{1-D mean-line design} -- thermodynamic and velocity-triangle
          analysis on the mean streamline;
    \item \textbf{Rotor-blade design} -- conversion of the flow angles into a
          blade profile (chord, stagger, pitch, blade count);
    \item \textbf{CFD evaluation} -- meshing and RANS solution of the resulting
          cascade to verify the loading and exit flow.
\end{itemize}

\begin{wrapfigure}{r}{0.55\textwidth}
    \centering
    \includegraphics[width=\linewidth]{engine.png}
    \caption{Industrial gas turbine of the SGT5-4000F class (cutaway).}
    \label{fig:engine}
\end{wrapfigure}

These steps are not executed once and in isolation: the design coefficients
(stage loading $\psi$, flow coefficient $\phi$ and degree of reaction $r$) are
iterated until the resulting stage satisfies a set of aerodynamic and mechanical
design guidelines. This report focuses on the \textbf{first stage} of a
four-stage high-pressure turbine sized for the conditions of an F-class
industrial gas turbine (representative of the Siemens SGT5-4000F), delivering
\SI{330}{MW} of shaft power at a mass flow of \SI{750}{kg/s}.




The whole design chain has been implemented in Python (the accompanying
notebook): the mean-line solver, the ParaBlade-based blade
construction~\cite{parablade} and the GMSH~\cite{gmsh}\,+\,SU2~\cite{su2} cascade
CFD. All numerical results, tables and figures in this report are produced
directly by that code.
